
\documentclass[UTF8,11pt,a4paper]{ctexart}

\usepackage{amsmath, amssymb, amsfonts}
\usepackage{geometry}
\usepackage{booktabs}
\usepackage{array}
\usepackage{longtable}
\usepackage{hyperref}
\usepackage{xcolor}
\usepackage{listings}
\usepackage{enumitem}
\usepackage{caption}
\usepackage{makecell}

\newcommand{\logSlots}{\mathsf{logSlots}}
\geometry{margin=1in}
\hypersetup{
    colorlinks=true,
    linkcolor=blue!60!black,
    urlcolor=blue!60!black,
    citecolor=blue!60!black
}

\lstset{
  basicstyle=\ttfamily\small,
  breaklines=true,
  frame=single,
  columns=fullflexible,
  keepspaces=true
}

\title{CKKS 参数选择教学文档：从 $N,n,Q,P,\Delta,L$ 到安全版 no-$B_k$ 实验}
\author{整理稿}
\date{\today}

\begin{document}
\maketitle

\begin{abstract}
本文档用于解释 CKKS 与 bootstrapping 实验中最常见的参数：
环维度 $N$、拆分后的 leaf 维度 $n$、拆分因子 $k$、密文模数链 $Q$、key-switching 辅助模数 $P$、最大公开模数 $QP$、scale $\Delta$、level 数 $L$、slot 数以及 bit security 之间的关系。
本文档的目标不是给出最终密码学安全证明，而是建立一套工程上可用的参数选择直觉，并给出当前 no-$B_k$ coefficient-split bootstrapping 实验应如何选安全参数。
\end{abstract}

\tableofcontents
\newpage

\section{最重要的结论}

先给出最实用的结论。

\begin{center}
\begin{tabular}{ll}
\toprule
问题 & 结论 \\
\midrule
$N$ 是什么？ & CKKS 的 ring degree，通常 $N=2^{\log N}$。\\
slot 数是多少？ & full complex CKKS 通常有 $N/2$ 个 complex slots。\\
$N$ 越大是不是消息更多？ & 是，但更重要的是安全预算更大，能承受更大的 $Q$ 或 $QP$。\\
$Q$ 是什么？ & ciphertext modulus chain 的乘积，支持 depth、precision 和 bootstrapping。\\
$P$ 是什么？ & key switching / relinearization / rotation keys 的 auxiliary modulus。\\
安全估计看 $Q$ 还是 $QP$？ & 通常保守看攻击者可见的最大公开模数，bootstrapping keys 常看 $QP$。\\
拆到 leaf 后看谁的安全？ & 看 leaf ring $R_n$ 上的 $(n,q_{\max})$，不能自动继承 top $N$。\\
$n=N/k$ 安全预算如何变？ & 粗略地，固定 128-bit 时允许的 $\log q$ 约随 $n$ 线性增长。\\
当前最推荐安全实验 & $\log N_{\text{top}}=17,\ k=2,\ \log n_{\text{leaf}}=16,\ \log(QP)\approx 1520$。\\
更大 $k=4$ 安全实验 & $\log N_{\text{top}}=18,\ k=4,\ \log n_{\text{leaf}}=16$，较重。\\
当前 $\log N=16,k=4,n=2^{14}$ & correctness prototype，不应作为 128-bit 安全实例。\\
\bottomrule
\end{tabular}
\end{center}

\section{CKKS 的几个核心参数}

\subsection{$N$：ring degree}

CKKS 通常工作在 2-power cyclotomic / negacyclic ring：

\[
R_Q = \mathbb{Z}_Q[X]/(X^N+1),
\qquad N=2^{\log N}.
\]

在 full complex CKKS 中，一个 ciphertext 最多可以编码：

\[
\frac{N}{2}
\]

个 complex slots。

\begin{center}
\begin{tabular}{rrrr}
\toprule
$\log N$ & $N$ & complex slots $N/2$ & 常见用途 \\
\midrule
13 & 8192 & 4096 & debug, toy correctness \\
14 & 16384 & 8192 & prototype, 小 leaf, 低/中等模数 \\
15 & 32768 & 16384 & 中等 CKKS，轻量 bootstrapping 或实验 leaf \\
16 & 65536 & 32768 & 常见 CKKS bootstrapping 安全维度 \\
17 & 131072 & 65536 & 大 batch；安全 $k=2$ no-$B_k$ top \\
18 & 262144 & 131072 & 更大 batch；安全 $k=4$ no-$B_k$ top \\
\bottomrule
\end{tabular}
\end{center}

\paragraph{核心直觉。}
$N$ 不是单纯的容量参数，它同时是：

\[
\boxed{
\text{容量参数} + \text{安全参数} + \text{性能成本参数}.
}
\]

$N$ 越大：
\begin{itemize}[leftmargin=2em]
    \item slots 越多；
    \item 可以承受更大的 ciphertext modulus；
    \item 可以支持更深的计算或 bootstrapping；
    \item 但 ciphertext、key、NTT、rotation、key-switching 和内存成本都变大。
\end{itemize}

\subsection{$n$ 与 $k$：拆分后的 leaf ring}

你的 no-$B_k$ coefficient-split 路线是：

\[
N \longrightarrow n=\frac{N}{k}.
\]

如果 $k=2^a$，则：

\[
\log n = \log N - a.
\]

例如：

\[
\log N=17,\quad k=2
\quad\Longrightarrow\quad
\log n=16.
\]

\[
\log N=18,\quad k=4
\quad\Longrightarrow\quad
\log n=16.
\]

这里最关键的是：

\[
\boxed{
\text{如果 leaf ciphertext / leaf keys 真正在 } R_n \text{ 上存在，则安全性要按 } n \text{ 估计。}
}
\]

不能说 top ring $N$ 安全，所以 leaf ring $n$ 也安全。

\subsection{$Q$：ciphertext modulus chain}

CKKS 中：

\[
Q = \prod_{i=0}^{L} q_i.
\]

通常每个 $q_i$ 是一个 RNS prime，例如 40-bit、50-bit 或 60-bit prime。

如果有 20 个 50-bit primes：

\[
\log_2 Q \approx 1000.
\]

$Q$ 的作用：
\begin{itemize}[leftmargin=2em]
    \item 支持更多 level；
    \item 支持更多乘法和 rescale；
    \item 支持 bootstrapping 中的 C2S、EvalMod、S2C；
    \item 增大 $Q$ 会降低 RLWE 安全性。
\end{itemize}

因此：

\[
\boxed{
Q \uparrow \Rightarrow \text{支持更深计算，但安全性下降，需要更大 } N.
}
\]

\subsection{$P$：key switching auxiliary modulus}

在 RNS CKKS 中，rotation key、relinearization key、bootstrapping key、ring switching key 等通常涉及一个 auxiliary modulus：

\[
P = \prod_j p_j.
\]

例如常见设置：

\[
\log_2 P \approx 4\times 61 \approx 244,
\]

或：

\[
\log_2 P \approx 5\times 61 \approx 305.
\]

安全估计时，经常不能只看 $Q$，而要看 evaluation keys 中出现的最大公开模数：

\[
q_{\max} \approx QP.
\]

这也是你 no-$B_k$ 分解安全分析中最重要的一点：

\[
\boxed{
\text{leaf bootstrapping keys 如果在 } R_{n,QP} \text{ 上生成，则 leaf 安全要看 } (n,QP).
}
\]

\subsection{$\Delta$：CKKS scale}

CKKS 把 real/complex message 乘上 scale $\Delta$ 编码成近似整数：

\[
m \mapsto \lfloor \Delta m \rceil.
\]

常见：

\[
\Delta=2^{30},\quad 2^{35},\quad 2^{40}.
\]

scale 越大：
\begin{itemize}[leftmargin=2em]
    \item 初始精度越高；
    \item 每次 multiplication 后 scale 变为 $\Delta^2$；
    \item 需要 rescale；
    \item 每次 rescale 大约消耗一个与 $\log\Delta$ 类似大小的 prime。
\end{itemize}

所以：

\[
\Delta \uparrow \Rightarrow Q \uparrow \Rightarrow N \uparrow.
\]

\subsection{$L$：level / multiplicative depth}

如果一个 CKKS computation 需要 $L$ 次 rescale，则 modulus chain 至少需要对应数量的 primes。

粗略地，如果每层 prime 约 $\log\Delta$ bits，那么：

\[
\log_2 Q \approx \log q_0 + L\log\Delta + \text{special/residual primes}.
\]

Bootstrapping 本身也需要很多 level，例如 C2S、EvalMod、S2C 都会消耗 levels。因此 CKKS bootstrapping 常见的 $\log(QP)$ 会到 1500 bits 左右。

\section{安全性依赖哪些数字？}

CKKS/RLWE 安全性主要看：

\[
\boxed{
N,\quad q_{\max},\quad \chi_s,\quad \chi_e,\quad \sigma,\quad \text{public samples / keys}.
}
\]

其中最关键的工程量是：

\[
\boxed{
N \text{ 和 } \log_2 q_{\max}.
}
\]

在有 evaluation keys 的 bootstrapping 系统中，保守地：

\[
q_{\max}
=
\max\{
Q_{\text{ciphertext}},
Q_{\text{public/eval keys}},
QP_{\text{switching keys}},
QP_{\text{bootstrapping keys}},
QP_{\text{split/merge keys}}
\}.
\]

通常可以先取：

\[
q_{\max}=QP.
\]

\subsection{为什么 $q$ 越大安全性越低？}

RLWE sample 可抽象为：

\[
(a,b)=(a,as+e)\pmod q.
\]

噪声 $e$ 小，模数 $q$ 大，意味着相对噪声比例变小：

\[
\frac{\|e\|}{q} \downarrow.
\]

攻击者更容易区分或恢复 secret。因此：

\[
\boxed{
q \uparrow \Rightarrow \text{安全性下降}.
}
\]

而 $N$ 越大，格维度越高，攻击越困难：

\[
\boxed{
N \uparrow \Rightarrow \text{安全性上升}.
}
\]

\section{128-bit 安全的参考表}

FHE Security Guidelines 给出了一张非常实用的表：在 Gaussian error standard deviation $\sigma=3.19$、uniform ternary secret 或 Gaussian secret 下，不同维度 $N$ 对应 Category 128/192/256 的最大 $\log_2(q)$。这里摘取最常用的 128-bit / ternary 行：

\begin{center}
\begin{tabular}{cc}
\toprule
Ring degree $N$ & 128-bit 下最大 $\log_2(q)$，ternary secret \\
\midrule
$2^{10}=1024$ & 26 \\
$2^{11}=2048$ & 53 \\
$2^{12}=4096$ & 106 \\
$2^{13}=8192$ & 214 \\
$2^{14}=16384$ & 430 \\
$2^{15}=32768$ & 868 \\
$2^{16}=65536$ & 1747 \\
$2^{17}=131072$ & 3523 \\
\bottomrule
\end{tabular}
\end{center}

\paragraph{注意。}
这张表不是你的最终安全证明。最终仍要跑 lattice estimator。但它是非常好的工程直觉表。

\subsection{最重要的近似规律}

从表中可以看到，固定 128-bit 安全时：

\[
\boxed{
\log_2 q_{\max}(N)
\text{ 大致随 } N \text{ 线性增长}.
}
\]

因此如果：

\[
n=\frac{N}{k},
\]

那么粗略地：

\[
\log_2 q_{\max}(n)
\approx
\frac{1}{k}\log_2 q_{\max}(N).
\]

这就是为什么：

\[
\boxed{
\text{拆到子环后，如果 leaf 仍用 top 的大 } QP,\text{ leaf 安全性通常会下降。}
}
\]

\section{把安全表套到你的 no-$B_k$ 实验}

\subsection{当前 correctness prototype：$\log N=16,k=4,\log n=14$}

参数：

\[
\log N_{\text{top}}=16,\quad k=4,\quad \log n_{\text{leaf}}=14.
\]

如果 leaf 使用：

\[
\log_2(QP)\approx 1520,
\]

但 $n=2^{14}$ 的 128-bit 参考预算只有：

\[
\log_2 q_{\max}\approx 430.
\]

因此：

\[
1520 \gg 430.
\]

所以这组应该定位为：

\[
\boxed{
\text{correctness / performance prototype，而不是 128-bit secure instantiation.}
}
\]

\subsection{安全版 $k=2$：$\log N=17,k=2,\log n=16$}

参数：

\[
\log N_{\text{top}}=17,\quad k=2,\quad \log n_{\text{leaf}}=16.
\]

如果：

\[
\log_2(QP)\approx 1520,
\]

而 $n=2^{16}$ 的 128-bit 预算约为：

\[
1747.
\]

因此：

\[
1520 < 1747.
\]

这就是目前最推荐的安全感更强的实验设置：

\[
\boxed{
\log N_{\text{top}}=17,\quad k=2,\quad \log n_{\text{leaf}}=16,\quad \log(QP)\approx 1520.
}
\]

你的实验中这组的 leaf boot params 是：

\[
\log n=16,\quad \logSlots=15,\quad \log QP=1520.0.
\]

同时 top 和 leaf Q/P moduli 完全一致，ring-basis compatibility check 通过。实验还显示 RAW normalization 对齐，即 $\gamma_{\mathrm{impl}}\approx 1$。

\subsection{安全版 $k=4$：$\log N=18,k=4,\log n=16$}

如果希望保留 $k=4$ 的更高并行度，同时 leaf 仍然使用安全常见的 $\log n=16$：

\[
k=4,\quad \log n=16
\quad\Rightarrow\quad
\log N_{\text{top}}=18.
\]

这是：

\[
\boxed{
\log N_{\text{top}}=18,\quad k=4,\quad \log n_{\text{leaf}}=16.
}
\]

优点：
\begin{itemize}[leftmargin=2em]
    \item leaf 可以承受 $\log(QP)\approx 1500$ 级别的 bootstrapping；
    \item $k=4$ 并行潜力更大；
    \item slots 更多，适合大 batch / high-throughput 场景。
\end{itemize}

缺点：
\begin{itemize}[leftmargin=2em]
    \item direct baseline 会非常慢；
    \item memory、evaluation keys、rotation keys 都会大幅增加；
    \item 实验可能很重。
\end{itemize}

\section{推荐参数表}

\begin{center}
\begin{tabular}{lrrrrl}
\toprule
名称 & top $\log N$ & $k$ & leaf $\log n$ & $\log(QP)$ & 定位 \\
\midrule
P0 prototype & 16 & 4 & 14 & 1520 & correctness only，不安全 \\
P1 small demo & 16 & 2 & 15 & 1520 & correctness/timing，不安全或低安全 \\
S1 recommended & 17 & 2 & 16 & 1520--1555 & 推荐安全版 $k=2$ \\
S2 heavy & 18 & 4 & 16 & 1520--1555 & 推荐安全版 $k=4$，很重 \\
S3 high capacity & 18 & 4 & 16 & 1730--1745 & 接近 $2^{16}$ leaf 128-bit 上界，需 estimator \\
\bottomrule
\end{tabular}
\end{center}

\paragraph{最推荐先跑的参数。}

\[
\boxed{
\log N_{\mathrm{top}}=17,\quad
k=2,\quad
\log n_{\mathrm{leaf}}=16,\quad
\log_2(QP)\approx 1520.
}
\]

这是最均衡的选择：leaf 维度足够大，correctness 已经跑通，speedup 明显，实验成本还不至于像 $\log N=18$ 那么重。

\section{为什么 $\log N=17,k=2$ 不应该和 $\log N=16$ 直接比 latency？}

$\log N=17$ 的 slot 数是：

\[
2^{16}=65536.
\]

$\log N=16$ 的 slot 数是：

\[
2^{15}=32768.
\]

所以：

\[
\log N=17
\]

处理的是两倍容量。直接比较：

\[
T(\log N=17,k=2)
\quad\text{vs}\quad
T(\log N=16)
\]

是不公平的。更合理的是比较：

\[
T_{\mathrm{factored}}(17,k=2)
\quad\text{vs}\quad
T_{\mathrm{direct}}(17),
\]

以及 throughput 角度比较：

\[
T_{\mathrm{factored}}(17,k=2)
\quad\text{vs}\quad
2T_{\mathrm{direct}}(16).
\]

\section{当前实验数据摘要}

\subsection{$\log N=17,k=2,\log n=16$}

\begin{center}
\begin{tabular}{lr}
\toprule
指标 & 数值 \\
\midrule
top $\log N$ & 17 \\
leaf $\log n$ & 16 \\
$k$ & 2 \\
$\log(QP)$ & 1520.0 \\
Ring basis compatibility & PASS \\
Best C2S alignment error & $\approx 4.0\times 10^{-11}$ \\
Best after-EvalMod alignment error & $\approx 2.0\times 10^{-5}$ \\
Direct final max error vs plaintext & $\approx 7\times 10^{-8}$ \\
Factored RAW max error vs plaintext & $\approx 7\times 10^{-7}$ \\
Direct vs factored RAW & $\approx 7\times 10^{-7}$ \\
Direct vs factored DIV-2 & $\approx 5.2\times 10^{-1}$ \\
Observed normalization & RAW is best, $\gamma\approx 1$ \\
Direct output level & 10 \\
Factored output level & 10 \\
Direct full time & $\approx 117.2$s \\
Factored full time & $\approx 23.3$s \\
Speedup & $\approx 5.03\times$ \\
\bottomrule
\end{tabular}
\end{center}

\subsection{$\log N=16,k=2,\log n=15$}

\begin{center}
\begin{tabular}{lr}
\toprule
指标 & 数值 \\
\midrule
top $\log N$ & 16 \\
leaf $\log n$ & 15 \\
$k$ & 2 \\
$\log(QP)$ & 1520.0 \\
Ring basis compatibility & PASS \\
Best C2S alignment error & $\approx 1.8\times 10^{-11}$ \\
Best after-EvalMod alignment error & $\approx 9.4\times 10^{-6}$ \\
Direct final max error vs plaintext & $\approx 3.2\times 10^{-8}$ \\
Factored RAW max error vs plaintext & $\approx 1.05$ \\
Factored DIV-2 max error vs plaintext & $\approx 2.4\times 10^{-7}$ \\
Direct vs factored DIV-2 & $\approx 2.4\times 10^{-7}$ \\
Observed normalization & DIV-2 is best, $\gamma\approx 2$ \\
Direct output level & 10 \\
Factored output level & 10 \\
Direct full time & $\approx 26.3$s \\
Factored full time & $\approx 7.22$s \\
Speedup & $\approx 3.64\times$ \\
\bottomrule
\end{tabular}
\end{center}

\section{参数选择流程：从应用需求到 $N,Q,n$}

\subsection{Step 1：确定应用需求}

先回答：

\begin{enumerate}[leftmargin=2em]
    \item 需要多少 slots？
    \item 需要多少 precision？
    \item 需要多少 multiplication depth？
    \item 是否需要 bootstrapping？
    \item 是追求 latency 还是 throughput？
\end{enumerate}

如果是 encrypted neural network inference / training，通常更偏 throughput，需要大量 slots 和 bootstrapping。

\subsection{Step 2：确定 scale $\Delta$}

常见选择：

\[
\Delta=2^{35}\quad\text{或}\quad 2^{40}.
\]

如果希望高精度，选 $2^{40}$；如果希望省 modulus，选 $2^{35}$ 或更小。

\subsection{Step 3：估计 $Q$}

如果有 $L$ 层乘法，每层 prime 大约接近 $\log\Delta$，则：

\[
\log Q
\approx
\log q_0 + L\log\Delta + \text{bootstrapping primes}.
\]

bootstrapping 参数中还要考虑 C2S、EvalMod、S2C、ModUp 等消耗。

\subsection{Step 4：加上 $P$ 得到 $QP$}

如果使用 key switching auxiliary primes：

\[
\log P \approx 4\times61=244
\quad \text{或}\quad
5\times61=305.
\]

最终要估：

\[
\log(QP)=\log Q+\log P.
\]

\subsection{Step 5：根据 $q_{\max}$ 选择 leaf $n$}

这是你的方案最重要的一步。

如果：

\[
\log(QP)\approx 1520,
\]

查表得到：

\[
2^{15}\text{ 只支持约 }868,
\]

\[
2^{16}\text{ 支持约 }1747.
\]

所以 leaf 至少应该选：

\[
\boxed{
\log n=16.
}
\]

\subsection{Step 6：根据 $k$ 得到 top $N$}

如果 leaf 固定为 $\log n=16$：

\[
k=2 \Rightarrow \log N=17,
\]

\[
k=4 \Rightarrow \log N=18.
\]

\subsection{Step 7：跑 estimator}

最终论文不能只靠表。应报告：

\[
\lambda_{\mathrm{scheme}}
=
\min\{
\lambda_{\mathrm{top}},
\lambda_{\mathrm{leaf}},
\lambda_{\mathrm{split}},
\lambda_{\mathrm{merge}},
\lambda_{\mathrm{evalkeys}}
\}.
\]

并列出：

\begin{center}
\begin{tabular}{lrrrrl}
\toprule
component & ring degree & $\log q_{\max}$ & secret & $\sigma$ & estimated bits \\
\midrule
top ciphertext/key & $2^{17}$ & 1520 & ternary & 3.19 & TBD \\
leaf BTS key & $2^{16}$ & 1520 & ternary & 3.19 & TBD \\
split key & ? & ? & ? & ? & TBD \\
merge key & ? & ? & ? & ? & TBD \\
\bottomrule
\end{tabular}
\end{center}

\section{no-$B_k$ 分解中的安全性表达}

你的方案可以写作：

\[
\mathrm{BTS}^{\mathrm{fac}}_N
=
\mathrm{Merge}_k
\circ
\left(
\bigoplus_{r=0}^{k-1}\mathrm{BTS}_n^{(r)}
\right)
\circ
\mathrm{Split}_k.
\]

其中：

\[
n=N/k.
\]

安全性不是只看 top ring：

\[
\mathrm{RLWE}_{N,Q}.
\]

还必须看 leaf ring：

\[
\mathrm{RLWE}_{n,Q_{\mathrm{leaf}}P_{\mathrm{leaf}}}.
\]

因此：

\[
\boxed{
\lambda_{\mathrm{scheme}}
=
\min_i \lambda_i.
}
\]

如果任何 leaf 或 switching key 只有 60-bit，那么整体就是 60-bit，不是 128-bit。

\section{参数和安全性的常见误区}

\subsection{误区 1：top $N$ 安全，所以 leaf $n$ 也安全}

错误。leaf 是小环实例，必须单独估。

\subsection{误区 2：只看输出 ciphertext level 的 $Q$}

错误。evaluation keys、rotation keys、bootstrapping keys 可能暴露更大的 $QP$ 模数，应保守看最大公开模数。

\subsection{误区 3：$N$ 越大只是消息更多}

不完整。$N$ 更重要的是提供安全预算来支持更大的 $Q$ 或 $QP$。

\subsection{误区 4：speedup 应该最多是 $k$}

不一定。direct bootstrapping 在大 $N$ 下可能超线性变慢，包含 C2S/S2C、rotations、key switching、NTT 和 memory movement。把一个大任务换成多个小 leaf 任务并行，speedup 可以超过 $k$。

\subsection{误区 5：normalization factor 永远是 $k$}

不是。你的实验已经说明：

\[
\gamma_{\mathrm{impl}}
=
\frac{
\mathrm{C2SScaling}_{leaf}
\mathrm{S2CScaling}_{leaf}
}{
\mathrm{C2SScaling}_{top}
\mathrm{S2CScaling}_{top}
}.
\]

对 $\log N=17,k=2$，$\gamma=1$；对 $\log N=16,k=2$，$\gamma=2$。这属于 implementation convention，不属于安全参数本身。

\section{推荐的实验命令}

\subsection{安全版 $k=2$}

\begin{lstlisting}[language=bash]
go run . \
  -logN=17 \
  -layers=1 \
  -leafOrder=bitrev \
  -parallelLeaves=true \
  -leafWorkers=2 \
  -reps=1 \
  -warmup=0
\end{lstlisting}

这对应：

\[
\log N_{\text{top}}=17,\quad k=2,\quad \log n_{\text{leaf}}=16.
\]

\subsection{安全版 $k=4$}

\begin{lstlisting}[language=bash]
go run . \
  -logN=18 \
  -layers=2 \
  -leafOrder=natural \
  -parallelLeaves=true \
  -leafWorkers=4 \
  -reps=1 \
  -warmup=0
\end{lstlisting}

这对应：

\[
\log N_{\text{top}}=18,\quad k=4,\quad \log n_{\text{leaf}}=16.
\]

这组会很重。

\section{Go 代码：打印参数和安全相关摘要}

建议在程序中打印：

\begin{lstlisting}[language=Go]
func printSecurityRelevantParams(name string, logN int, logQP float64, logQ float64, logP float64) {
    N := 1 << logN
    slots := N / 2
    fmt.Printf("\n=== %s security-relevant parameters ===\n", name)
    fmt.Printf("LogN      = %d\n", logN)
    fmt.Printf("N         = %d\n", N)
    fmt.Printf("Slots     = %d\n", slots)
    fmt.Printf("logQ      = %.1f\n", logQ)
    fmt.Printf("logP      = %.1f\n", logP)
    fmt.Printf("logQP     = %.1f\n", logQP)
}
\end{lstlisting}

对你的实验，至少打印：

\begin{lstlisting}[language=Go]
printSecurityRelevantParams("TOP", topLogN, paramsTop.LogQP(), logQTop, logPTop)
printSecurityRelevantParams("LEAF", leafLogN, paramsLeaf.LogQP(), logQLeaf, logPLeaf)
\end{lstlisting}

并在 summary 中给出：

\[
\log N_{\text{top}},
\quad
\log n_{\text{leaf}},
\quad
k,
\quad
\log Q,
\quad
\log P,
\quad
\log(QP).
\]

\section{最终建议}

当前阶段建议你把实验分成两类：

\subsection{Correctness / mechanism prototype}

这些参数可以小一些：

\[
\log N=16,\quad k=4,\quad \log n=14.
\]

用途：
\begin{itemize}[leftmargin=2em]
    \item 验证 no-$B_k$ 交换图；
    \item 验证 C2S/EvalMod/S2C 拆分正确性；
    \item 验证 normalization trace；
    \item 做快速 timing。
\end{itemize}

但不声称 128-bit security。

\subsection{Security-aware prototype}

推荐：

\[
\boxed{
\log N=17,\quad k=2,\quad \log n=16,\quad \log(QP)\approx1520.
}
\]

用途：
\begin{itemize}[leftmargin=2em]
    \item 展示更接近 128-bit 参数的 no-$B_k$ 分解；
    \item leaf ring 维度足够承受 bootstrapping 级别模数；
    \item 和 direct $\log N=17$ 做公平算法对比；
    \item 报告 output level、precision、speedup。
\end{itemize}

\subsection{Full secure high-throughput target}

推荐：

\[
\boxed{
\log N=18,\quad k=4,\quad \log n=16.
}
\]

用途：
\begin{itemize}[leftmargin=2em]
    \item 展示高吞吐 / 大 batch 场景；
    \item leaf 仍为安全常见维度；
    \item $k=4$ 并行潜力更大。
\end{itemize}

缺点是计算和内存会很重。

\section{一句话总结}

\[
\boxed{
N \text{ 是安全预算和容量预算，} QP \text{ 是消耗安全预算，} n=N/k \text{ 后必须重新估 leaf 安全。}
}
\]

对当前 no-$B_k$ bootstrapping：

\[
\boxed{
\log(QP)\approx 1520 \text{ 时，leaf 最好至少取 } \log n=16.
}
\]

因此：

\[
\boxed{
k=2 \Rightarrow \log N_{\text{top}}=17,
\qquad
k=4 \Rightarrow \log N_{\text{top}}=18.
}
\]

\section*{参考资料}

\begin{enumerate}[leftmargin=2em]
    \item HomomorphicEncryption.org, \emph{Homomorphic Encryption Standard}. \url{https://homomorphicencryption.org/standard/}
    \item Albrecht et al., \emph{Security Guidelines for Implementing Homomorphic Encryption}, ePrint 2024/463. Table 5.2 gives maximal $\log_2(q)$ for dimensions and security categories. \url{https://eprint.iacr.org/2024/463}
    \item Lattigo v6 bootstrapping package documentation. It states that bootstrapping parameters have independent cryptographic parameters and users must ensure target security. \url{https://pkg.go.dev/github.com/tuneinsight/lattigo/v6/circuits/ckks/bootstrapping}
    \item 当前实验日志：\texttt{dftparams\_logN17\_k2.txt}, \texttt{dftparams\_logN16\_k2.txt}, \texttt{normtrace\_logN17\_k2.txt}, \texttt{normtrace\_logN16\_k2\_timing.txt}.
\end{enumerate}

\end{document}
